% Generated by `taskgen course examples` from factoring.pollard_pm1 (seed 7).
% Do not edit: change examples.yaml and regenerate.
\begin{example}[алгоритм $p - 1$ с показателями $j!$]\label{ex:pm1-factorial}
Найти нетривиальный делитель числа $N = 221$ с помощью $p-1$ алгоритма Полларда.

\par\noindent\textit{Решение.} Вычисляем $a_j = a_{j-1}^{\,j} \bmod N$ (т.е. $a^{j!}$) и $d_j = \gcd(a_j - 1, N)$, пока $d_j = 1$; при $1 < d_j < N$ найден делитель (для $p \mid N$ с $(p-1) \mid j!$ имеем $p \mid a^{j!} - 1$), при $d_j = N$ меняем основание.\par\noindent $a = 2$:\[\begin{array}{c|rrr}j & 2 & 3 & 4 \\ \hline a_j & 4 & 64 & 1 \\ d_j & 1 & 1 & 221\end{array}\]$d = N$ --- неудача, берём следующее основание.\par\noindent $a = 3$:\[\begin{array}{c|rr}j & 2 & 3 \\ \hline a_j & 9 & 66 \\ d_j & 1 & 13\end{array}\]Делитель $13$, $N / 13 = 17$.\[N = p \cdot q = 13 \cdot 17 = 221.\]
\end{example}
